\section{Relations and Equivalence}
\chapterline

In everyday life, we constantly decide that two different things are ``the same'' for a specific purpose. 

A one-dollar bill printed in 2017 and a crisp bill printed in 2024 are made of different atoms and carry different serial numbers. But if you buy a coffee, the cashier does not inspect the serial numbers. For the purpose of commerce, both bills are identical.

Similarly, on an analog clock, the numbers $1$, $13$, and $25$ represent entirely different quantities of hours. Yet when reading the hour hand, all three point to the exact same position on the dial. The differences of $12$ hours have been deliberately collapsed.

In mathematics, the tool that formalizes this act of purposeful simplification is the \textbf{equivalence relation}.

\begin{historybox}
The conceptual origin of equivalence relations traces directly to a young Carl Friedrich Gauss. In his 1801 masterpiece, \emph{Disquisitiones Arithmeticae}, the 24-year-old Gauss revolutionized number theory by introducing modular arithmetic and the congruence symbol $a \equiv b \pmod n$. 

Before Gauss, mathematicians treated remainders as awkward numerical byproducts of Euclidean division. When dividing $17$ by $5$, the remainder $2$ was merely what was left over. Gauss made an extraordinary conceptual leap: numbers that leave the same remainder modulo $n$ should not be treated as separate integers with an incidental shared feature. They should be recognized as members of a single, unified mathematical organism: an \emph{equivalence class}.

Contemporaries initially balked at Gauss's new notation, asking why one could not simply write the equation $a - b = k n$. Gauss understood that the notation $\equiv$ reflected a profound structural reality: congruence behaves exactly like equality. It allows one to perform arithmetic on the classes themselves: $[a] + [b] = [a+b]$ and $[a][b] = [ab]$. 

Over the next century, mathematicians like Richard Dedekind, Henri Poincar\'e, and Emmy Noether realized that Gauss's construction was not unique to integers. It was the universal blueprint for manufacturing new spaces in geometry and abstract algebra: whenever you want to ignore extraneous details, factor out a symmetry, or glue geometric boundaries together, you define an equivalence relation and pass to the quotient.
\end{historybox}

\subsection{Relations in General}

Before studying equivalence, we need the general concept of a relation between elements of a set.

\begin{definition}[Relation]
A \textbf{binary relation} $R$ on a set $A$ is a subset of the Cartesian product:
\[
R \subseteq A \times A.
\]
If $(a, b) \in R$, we write $a \sim b$ or $a \, R \, b$, and say that ``$a$ is related to $b$.''
\end{definition}

A relation simply declares which pairs of elements are linked. Examples include:
\begin{itemize}
\item Arithmetic orders: $x < y$ or $x \le y$ on $\R$.
\item Divisibility: $a \mid b$ on the integers $\Z$.
\item Geometric connectivity: two points in a space connected by a continuous path.
\item Set containment: $S \subseteq T$ on the power set $\pow(U)$.
\end{itemize}

\subsection{The Three Axioms of Equivalence}

To behave like an idealized equality, a relation must satisfy three essential structural axioms.

\begin{definition}[Equivalence Relation]
A relation $\sim$ on a set $A$ is an \textbf{equivalence relation} if it satisfies:
\begin{enumerate}[label=(\roman*)]
\item \textbf{Reflexivity:} For all $a \in A$,
\[
a \sim a.
\]
(Every object is related to itself.)
\item \textbf{Symmetry:} For all $a, b \in A$,
\[
a \sim b \implies b \sim a.
\]
(The relation is two-way; perspective does not matter.)
\item \textbf{Transitivity:} For all $a, b, c \in A$,
\[
(a \sim b \text{ and } b \sim c) \implies a \sim c.
\]
(Chains of relation do not break.)
\end{enumerate}
\end{definition}

\begin{intuitionbox}
Think of reflexivity as self-recognition, symmetry as mutual agreement, and transitivity as contagion. Transitivity is the engine that prevents the world from fragmenting: if $a$ belongs to the same club as $b$, and $b$ belongs to the same club as $c$, then $a$ and $c$ must share membership in that club.
\end{intuitionbox}

\begin{trapbox}
\textbf{The Trap of Naive Similarity:}
It is tempting to think that any intuitive notion of ``being similar'' defines an equivalence relation. Consider the relation on $\R$ defined by $|x - y| \le 1$ (``$x$ and $y$ are within distance 1 of each other''). 
Reflexivity holds ($|x - x| = 0 \le 1$), and symmetry holds ($|y - x| = |x - y| \le 1$). 
Yet transitivity fails completely: $0 \sim 0.8$ and $0.8 \sim 1.5$, but $|0 - 1.5| = 1.5 > 1$, so $0 \not\sim 1.5$!
Without transitivity, elements cannot be partitioned into disjoint, non-overlapping islands. Similarity bleeds across boundaries; only true equivalence builds clean quotient spaces.
\end{trapbox}

\subsection{Equivalence Classes and Quotients}

Once you have an equivalence relation on $A$, the set breaks apart into clusters called \emph{equivalence classes}.

\begin{definition}[Equivalence Class]
Let $\sim$ be an equivalence relation on $A$. For any $a \in A$, the \textbf{equivalence class} of $a$, denoted $[a]$ or $\bar{a}$, is the set of all elements related to $a$:
\[
[a] = \{x \in A : x \sim a\}.
\]
The element $a$ is called a \textbf{representative} of the class $[a]$.
\end{definition}

\begin{center}
\begin{tikzpicture}[scale=0.92, >=Latex]
  % Stage 1: Partitioned Space X
  \begin{scope}[yshift=0cm]
    \node[anchor=west, font=\small\bfseries] at (-0.5, 3.2) {Stage 1: Equivalence Classes Partition the Space $X$};
    \draw[rounded corners=4pt, fill=gray!3, draw=gray!25] (-0.6, -0.6) rectangle (12.2, 3.0);
    
    % Four cells / tiles
    % Tile 1: [a]
    \draw[thick, sourceblue, fill=sourceblue!15, rounded corners=3pt] (0.2, 0.1) rectangle (2.6, 2.3);
    \node[sourceblue!90!black, font=\scriptsize\bfseries] at (1.4, 2.0) {Class $[a]$};
    \node[circle, fill=sourceblue, inner sep=1.5pt, label={below:\tiny $a$}] (a1) at (0.8, 0.7) {};
    \node[circle, fill=sourceblue, inner sep=1.5pt, label={right:\tiny $x$}] (a2) at (1.9, 1.2) {};
    \draw[thick, sourceblue!70, dashed] (a1) -- node[above left, font=\tiny] {$\sim$} (a2);
    
    % Tile 2: [b]
    \draw[thick, targetgreen, fill=targetgreen!15, rounded corners=3pt] (3.2, 0.1) rectangle (5.6, 2.3);
    \node[targetgreen!90!black, font=\scriptsize\bfseries] at (4.4, 2.0) {Class $[b]$};
    \node[circle, fill=targetgreen, inner sep=1.5pt, label={below:\tiny $b$}] (b1) at (3.8, 0.8) {};
    \node[circle, fill=targetgreen, inner sep=1.5pt, label={above:\tiny $y$}] (b2) at (4.8, 1.4) {};
    \node[circle, fill=targetgreen, inner sep=1.5pt, label={below:\tiny $z$}] (b3) at (5.0, 0.5) {};
    \draw[thick, targetgreen!70, dashed] (b1) -- (b2) -- (b3) -- (b1);
    
    % Tile 3: [c]
    \draw[thick, quotientorange, fill=quotientorange!15, rounded corners=3pt] (6.2, 0.1) rectangle (8.6, 2.3);
    \node[quotientorange!90!black, font=\scriptsize\bfseries] at (7.4, 2.0) {Class $[c]$};
    \node[circle, fill=quotientorange, inner sep=1.5pt, label={below:\tiny $c$}] (c1) at (7.0, 0.9) {};
    \node[circle, fill=quotientorange, inner sep=1.5pt, label={above:\tiny $w$}] (c2) at (7.8, 1.3) {};
    \draw[thick, quotientorange!70, dashed] (c1) -- node[below right, font=\tiny] {$\sim$} (c2);
    
    % Tile 4: [d]
    \draw[thick, structurepurple, fill=structurepurple!15, rounded corners=3pt] (9.2, 0.1) rectangle (11.6, 2.3);
    \node[structurepurple!90!black, font=\scriptsize\bfseries] at (10.4, 2.0) {Class $[d]$};
    \node[circle, fill=structurepurple, inner sep=1.5pt, label={below:\tiny $d$}] (d1) at (10.4, 1.0) {};
  \end{scope}

  % Middle Corridor: Canonical Projection Map
  \begin{scope}[yshift=-1.4cm]
    \draw[->, very thick, gray!70] (1.4, 0.6) -- node[right, font=\scriptsize] {$\pi$} (1.4, -0.4);
    \draw[->, very thick, gray!70] (4.4, 0.6) -- node[right, font=\scriptsize] {$\pi$} (4.4, -0.4);
    \draw[->, very thick, gray!70] (7.4, 0.6) -- node[right, font=\scriptsize] {$\pi$} (7.4, -0.4);
    \draw[->, very thick, gray!70] (10.4, 0.6) -- node[right, font=\scriptsize] {$\pi$} (10.4, -0.4);
  \end{scope}

  % Stage 2: Quotient Space X / ~
  \begin{scope}[yshift=-3.3cm]
    \node[anchor=west, font=\small\bfseries] at (-0.5, 1.6) {Stage 2: The Quotient Space $X/\!\sim$ (Each Class Becomes a Point)};
    \draw[rounded corners=4pt, fill=gray!3, draw=gray!25] (-0.6, -0.7) rectangle (12.2, 1.4);
    
    \node[circle, draw=sourceblue, fill=sourceblue!30, inner sep=3pt, label={below:\footnotesize $[a]$}] at (1.4, 0.3) {};
    \node[circle, draw=targetgreen, fill=targetgreen!30, inner sep=3pt, label={below:\footnotesize $[b]$}] at (4.4, 0.3) {};
    \node[circle, draw=quotientorange, fill=quotientorange!30, inner sep=3pt, label={below:\footnotesize $[c]$}] at (7.4, 0.3) {};
    \node[circle, draw=structurepurple, fill=structurepurple!30, inner sep=3pt, label={below:\footnotesize $[d]$}] at (10.4, 0.3) {};
  \end{scope}
\end{tikzpicture}
\end{center}

\begin{theorem}[Fundamental Theorem of Equivalence Relations]
Let $\sim$ be an equivalence relation on $A$.
\begin{enumerate}[label=(\alph*)]
\item For any $a, b \in A$, $a \sim b \iff [a] = [b]$.
\item Any two equivalence classes are either strictly disjoint or completely identical:
\[
[a] \cap [b] \neq \emptyset \implies [a] = [b].
\]
\item The collection of all equivalence classes forms a \textbf{partition} of $A$: they are pairwise disjoint and their union is all of $A$.
\end{enumerate}
\end{theorem}

\begin{definition}[Quotient Set and Projection]
The collection of all equivalence classes of $A$ under $\sim$ is called the \textbf{quotient set}, denoted $A / \!\sim$:
\[
A / \!\sim \; = \{[a] : a \in A\}.
\]
The function $\pi: A \to A / \!\sim$ given by $\pi(a) = [a]$ is called the \textbf{canonical projection map}.
\end{definition}

In the quotient space, each entire equivalence class has collapsed into a single mathematical point.

\subsection{A Canonical Example: Modular Arithmetic}

Let $n$ be a positive integer. We say two integers $a, b \in \Z$ are \emph{congruent modulo $n$}, written $a \equiv b \pmod n$, if their difference is a multiple of $n$:
\[
a \equiv b \pmod n \iff n \mid (a - b).
\]

\begin{theorem}
Congruence modulo $n$ is an equivalence relation on $\Z$.
\end{theorem}

\begin{proof}
We check all three axioms systematically.

\textbf{1. Reflexivity:} Let $a \in \Z$. Then $a - a = 0 = n \cdot 0$. Because $0$ is an integer, $n \mid (a - a)$, so $a \equiv a \pmod n$.

\textbf{2. Symmetry:} Let $a, b \in \Z$ and assume $a \equiv b \pmod n$. 
By definition, $n \mid (a - b)$, meaning $a - b = n k$ for some integer $k$. 
Multiplying both sides by $-1$:
\[
b - a = -(a - b) = -nk = n(-k).
\]
Since $-k \in \Z$, $n \mid (b - a)$, which proves $b \equiv a \pmod n$.

\textbf{3. Transitivity:} Let $a, b, c \in \Z$, and assume $a \equiv b \pmod n$ and $b \equiv c \pmod n$.
Then $a - b = n j$ and $b - c = n k$ for integers $j, k \in \Z$. 
Adding the two equations:
\[
a - c = (a - b) + (b - c) = n j + n k = n(j + k).
\]
Since $j + k$ is an integer, $n \mid (a - c)$, confirming $a \equiv c \pmod n$.

Since reflexivity, symmetry, and transitivity all hold, congruence modulo $n$ is an equivalence relation.
\end{proof}

The quotient set $\Z / \equiv_n$ is the familiar ring $\Z_n = \{[0], [1], \dots, [n-1]\}$, containing exactly $n$ equivalence classes.

\subsection{Why Mathematicians Care: The Universal Construction}

Equivalence relations and quotient sets are how mathematicians manufacture new mathematical universes from old ones:
\begin{itemize}
\item \textbf{Constructing the Rational Numbers $\Q$:} Start with pairs of integers $(p, q)$ with $q \neq 0$. Declare $(p, q) \sim (r, s) \iff ps = qr$. The quotient set is precisely $\Q$.
\item \textbf{Topology and Manifolds:} Take a flat sheet of paper $[0, 1] \times [0, 1]$ and declare $(0, y) \sim (1, y)$ (glue left edge to right edge) and $(x, 0) \sim (x, 1)$ (glue bottom edge to top edge). The resulting quotient space is a 2-dimensional torus (a donut)!
\item \textbf{Abstract Algebra:} Quotients of groups by normal subgroups $G/N$ and rings by ideals $R/I$ are quotient sets equipped with lawful algebraic operations.
\end{itemize}

\subsection{What Quotients Unlocked, Their Limits, and Where We Go Next}

Equivalence relations and quotient sets unlocked one of the greatest creative engines in modern mathematics:
\begin{itemize}
\item \textbf{Constructing new numbers:} We construct the rational numbers $\Q$ as quotient pairs $(p, q) \sim (r, s) \iff ps = qr$, and the real numbers $\R$ as equivalence classes of Cauchy sequences that differ by a null sequence.
\item \textbf{Topological manifolds:} We build spheres, tori, Klein bottles, and projective spaces by declaring boundary points of flat polygons to be equivalent.
\item \textbf{Algebraic classification:} In abstract algebra (Chapter~21), the First Isomorphism Theorem ($G / \ker\phi \cong \operatorname{im}\phi$) demonstrates that every homomorphism is secretly a quotient projection followed by an isomorphism.
\end{itemize}

Yet equivalence relations possess an inescapable structural boundary: they are strictly \textbf{horizontal}. 

An equivalence relation takes a space and flattens it into peer-to-peer classes where every element inside a class has equal status. It is fundamentally incapable of modeling \textbf{vertical hierarchy}:
\begin{itemize}
\item Which elements contain, dominate, or precede other elements?
\item How do we model dependency graphs, compilers, or task scheduling?
\item How do we describe divisibility chains, where $2 \mid 4$ and $4 \mid 12$, but $12$ does not divide $4$?
\end{itemize}

To model vertical hierarchy, we must break the symmetry of equivalence. If $a$ precedes $b$, we cannot allow $b$ to precede $a$ unless they are the exact same object. In Chapter~16, we introduce \textbf{Orders: Another Kind of Relation}, trading symmetry for antisymmetry to explore posets, chains, and lattices.

\subsection{Exercises}

\begin{enumerate}
\item For each relation on $\R$, determine which of the three properties (reflexive, symmetric, transitive) hold:
\begin{enumerate}[label=(\alph*)]
\item $x \sim y \iff |x - y| \le 1$.
\item $x \sim y \iff x - y \in \Q$.
\item $x \sim y \iff xy > 0$.
\end{enumerate}
\item Let $f: A \to B$ be a function. Define a relation $\sim$ on $A$ by
\[
a_1 \sim a_2 \iff f(a_1) = f(a_2).
\]
Prove that $\sim$ is an equivalence relation. What do its equivalence classes represent geometrically?
\item Let $S$ be a set with 4 elements. How many distinct equivalence relations can be defined on $S$? (Hint: Count the number of partitions of a 4-element set, known as the Bell number $B_4$.)
\end{enumerate}

\newpage
